% © 2026 Ing. David Jaroš — CC BY-NC-ND 4.0
%
% This work is licensed under a Creative Commons Attribution-NonCommercial-NoDerivatives
% 4.0 International License (CC BY-NC-ND 4.0).
%
% License History: Earlier drafts (up to v0.3) were released under CC BY 4.0.
% From v0.4 onward, all material is released under CC BY-NC-ND 4.0 to protect
% the integrity of the theoretical work during ongoing academic development.
%
% See LICENSE.md for full license text.

\ifdefined\INCLUDEMODE
\else
  \documentclass[12pt,a4paper]{article}
  \usepackage[a4paper, margin=2.5cm]{geometry}
  \usepackage{amsmath, amssymb, amsthm}
  \usepackage{mathtools}
  \usepackage{hyperref}
  \usepackage{xcolor}
  \usepackage{mdframed}
  \usepackage{booktabs}

  \newtheorem{theorem}{Theorem}[section]
  \newtheorem{lemma}[theorem]{Lemma}
  \newtheorem{proposition}[theorem]{Proposition}
  \newtheorem{corollary}[theorem]{Corollary}
  \theoremstyle{definition}
  \newtheorem{definition}[theorem]{Definition}
  \theoremstyle{remark}
  \newtheorem{remark}[theorem]{Remark}

  \newmdenv[backgroundcolor=yellow!20, linecolor=orange, linewidth=1.5pt]{gapbox}
  \newmdenv[backgroundcolor=green!15, linecolor=green!60!black, linewidth=1.5pt]{resultbox}
  \newmdenv[backgroundcolor=blue!8, linecolor=blue!50, linewidth=1.5pt]{insightbox}

  \title{\textbf{Step 2: SU(3) from Octonion Restriction --- $\mathbb{C}\otimes\mathbb{H} \subset \mathbb{C}\otimes\mathbb{O}$}\\
         \large The Key Identity $SU(3) = \mathrm{Aut}(\mathbb{O})|_{\mathrm{fix\,}e_7}$}
  \author{David Jaro\v{s}}
  \date{2026-03-06}

  \begin{document}
  \maketitle

  \begin{abstract}
  We investigate the derivation of $SU(3)$ via octonion restriction:
  $\mathbb{C}\otimes\mathbb{H}$ is a subalgebra of $\mathbb{C}\otimes\mathbb{O}$
  (complex octonions), and the stabilizer of $e_7$ in $G_2 = \mathrm{Aut}(\mathbb{O})$
  is $SU(3)$, which acts on the complement $\mathbb{C}^3 \cong
  \mathrm{span}\{e_4, e_5, e_6\}$.  We work through the construction,
  identify the 8 generators with Gell-Mann matrices, and assess applicability
  to UBT.  Verdict: \textbf{SKETCH} --- the algebra works, but octonions
  are non-associative and outside the UBT foundation ($\mathbb{C}\otimes\mathbb{H}$);
  they serve as an auxiliary tool only.  The dynamical content of $SU(3)$
  gauge invariance requires the involution approach of
  \texttt{appendix\_G\_internal\_color\_symmetry.tex}.
  \end{abstract}
\fi

%──────────────────────────────────────────────────────────────────────────────
\section{Complex Octonions $\mathbb{C}\otimes\mathbb{O}$}
\label{sec:coct}
%──────────────────────────────────────────────────────────────────────────────

The octonions $\mathbb{O}$ form the unique normed division algebra of
real dimension 8.  Denote the basis of $\mathbb{O}$ as
$\{1, e_1, e_2, e_3, e_4, e_5, e_6, e_7\}$ where
$e_i^2 = -1$ for $i = 1,\ldots,7$.

The complex octonions are $\mathbb{C}\otimes\mathbb{O}$, with
\begin{equation}
  \dim_\mathbb{R}(\mathbb{C}\otimes\mathbb{O}) \;=\; 16,
  \qquad
  \dim_\mathbb{C}(\mathbb{C}\otimes\mathbb{O}) \;=\; 8.
\end{equation}

\begin{insightbox}
\textbf{Key dimensions:}
\begin{itemize}
  \item $\mathbb{C}\otimes\mathbb{O}$: 16 real / 8 complex dimensions.
  \item $\mathbb{C}\otimes\mathbb{H}$: 8 real / 4 complex dimensions.
  \item Complement: 8 real / 4 complex dimensions, spanned by
        $\{e_4, e_5, e_6, e_7\}$ (and their $i$-multiples).
\end{itemize}
\end{insightbox}

%──────────────────────────────────────────────────────────────────────────────
\section{The Subalgebra $\mathbb{C}\otimes\mathbb{H} \subset \mathbb{C}\otimes\mathbb{O}$}
\label{sec:subalgebra}
%──────────────────────────────────────────────────────────────────────────────

\begin{proposition}[Associativity of $\mathbb{C}\otimes\mathbb{H}$]
\label{prop:assoc}
The identification $\{1, e_1, e_2, e_3\} \leftrightarrow \{1, I, J, K\}$
embeds $\mathbb{H}$ as an associative subalgebra of $\mathbb{O}$.  Therefore
$\mathbb{C}\otimes\mathbb{H}$ is an associative subalgebra of $\mathbb{C}\otimes\mathbb{O}$.
\end{proposition}

\begin{proof}
The quaternionic subalgebra $\{1, e_1, e_2, e_3\} \subset \mathbb{O}$ is
closed under the octonion product: $e_1 e_2 = e_3$, $e_2 e_3 = e_1$,
$e_3 e_1 = e_2$ (from the standard Fano plane multiplication table).
Since any subalgebra of $\mathbb{O}$ generated by two elements is
associative (Artin's theorem), and $\mathbb{H} = \langle e_1, e_2\rangle$
requires only two generators, $\mathbb{H}$ is associative inside $\mathbb{O}$.
\end{proof}

The remaining basis elements $\{e_4, e_5, e_6, e_7\}$ span the
\emph{complement} of $\mathbb{H}$ in $\mathbb{O}$:
\begin{equation}
  \mathbb{O} \;=\; \mathbb{H} \;\oplus\; e_4\mathbb{H}
  \qquad (\text{Cayley--Dickson decomposition}).
\end{equation}

%──────────────────────────────────────────────────────────────────────────────
\section{The Group $G_2 = \mathrm{Aut}(\mathbb{O})$}
\label{sec:g2}
%──────────────────────────────────────────────────────────────────────────────

\begin{definition}[$G_2$ as automorphism group]
\begin{equation}
  G_2 \;:=\; \mathrm{Aut}(\mathbb{O})
       \;=\; \{g\in GL(\mathbb{O}) : g(xy) = g(x)g(y) \text{ for all } x,y\in\mathbb{O}\}.
\end{equation}
This is a compact simple Lie group with $\dim G_2 = 14$.
\end{definition}

Key properties:
\begin{itemize}
  \item Every $g\in G_2$ fixes $1 \in \mathbb{O}$ (automorphisms preserve the identity).
  \item $G_2 \subset SO(7)$ (acts on the 7 imaginary octonionic units by isometries).
  \item The Lie algebra $\mathfrak{g}_2$ has 14 generators.
\end{itemize}

%──────────────────────────────────────────────────────────────────────────────
\section{$SU(3)$ as Stabilizer in $G_2$}
\label{sec:su3stab}
%──────────────────────────────────────────────────────────────────────────────

\begin{theorem}[$SU(3) = \{g\in G_2 : g(e_7) = e_7\}$]
\label{thm:su3stab}
The stabilizer of $e_7$ in $G_2$ is isomorphic to $SU(3)$:
\begin{equation}
  SU(3) \;\cong\; \mathrm{Stab}_{G_2}(e_7)
         \;=\; \{g\in G_2 : g(e_7) = e_7\}.
\end{equation}
\end{theorem}

\begin{proof}[Proof sketch]
An automorphism $g\in G_2$ fixing $e_7$ must permute
$\{e_4, e_5, e_6, e_4e_7, e_5e_7, e_6e_7\}$ compatibly with the
octonion multiplication table.  The octonion multiplication restricted
to the 6-dimensional real space $\mathrm{span}\{e_4,\ldots,e_7,e_4e_7,\ldots\}$
gives the constraint that $g$ preserves the associative 3-form
$\phi \in \Lambda^3(\mathbb{R}^6)^*$.  The subgroup of $SO(6)$
preserving $\phi$ is exactly $SU(3)$.  See~\cite{Baez2002} for the
complete proof.
\end{proof}

%──────────────────────────────────────────────────────────────────────────────
\section{Action on the Complement $\mathbb{C}^3$}
\label{sec:action}
%──────────────────────────────────────────────────────────────────────────────

\begin{proposition}[Fundamental representation]
\label{prop:fund}
The stabilizer $SU(3) \cong \mathrm{Stab}_{G_2}(e_7)$ acts on
\begin{equation}
  \mathbb{C}^3 \;\cong\; \mathrm{span}_\mathbb{C}\{e_4, e_5, e_6\}
\end{equation}
in the fundamental representation $\mathbf{3}$.
\end{proposition}

\begin{proof}[Proof sketch]
An element $g\in SU(3) \subset G_2$ fixes $e_7$ and the quaternionic
subalgebra $\{1, e_1, e_2, e_3\}$.  It therefore acts on the complement
$\mathrm{span}_\mathbb{R}\{e_4, e_5, e_6, e_4e_7, e_5e_7, e_6e_7\}
\cong \mathbb{C}^3$ (where $e_7$ plays the role of $i$).
The resulting action is by $3\times 3$ unitary matrices with determinant 1,
i.e., in the fundamental $\mathbf{3}$ of $SU(3)$.
\end{proof}

%──────────────────────────────────────────────────────────────────────────────
\section{Identification with Gell-Mann Matrices}
\label{sec:gellmann}
%──────────────────────────────────────────────────────────────────────────────

The 8 generators of $\mathfrak{su}(3) \subset \mathfrak{g}_2$ acting on
$\mathbb{C}^3 \cong \mathrm{span}\{e_4, e_5, e_6\}$ can be written
in the basis $\{e_4, e_5, e_6\}$.  Using the octonion multiplication
table restricted to this complement:

\begin{align}
  T_a &= \tfrac{1}{2}\lambda_a, \quad a = 1,\ldots,8,
  \label{eq:generators}
\end{align}
where $\lambda_a$ are the Gell-Mann matrices (see
\texttt{step1\_generators\_from\_8D.tex}, eq.~(3)).  The identification
is:
\begin{center}
\begin{tabular}{cll}
  \toprule
  Generator $T_a$ & Octonion derivation & Gell-Mann matrix \\
  \midrule
  $T_1$ & $[L_{e_4}, L_{e_5}]|_{V_c}$ & $\lambda_1/2$ \\
  $T_2$ & $[L_{e_4}, R_{e_5}]|_{V_c}$ & $\lambda_2/2$ \\
  $T_3$ & $[L_{e_4}, L_{e_4}^\dagger]|_{V_c}$ & $\lambda_3/2$ \\
  $T_4$ & $[L_{e_4}, L_{e_6}]|_{V_c}$ & $\lambda_4/2$ \\
  $T_5$ & $[L_{e_4}, R_{e_6}]|_{V_c}$ & $\lambda_5/2$ \\
  $T_6$ & $[L_{e_5}, L_{e_6}]|_{V_c}$ & $\lambda_6/2$ \\
  $T_7$ & $[L_{e_5}, R_{e_6}]|_{V_c}$ & $\lambda_7/2$ \\
  $T_8$ & $[L_{e_4}, L_{e_4}^\dagger] + [L_{e_5}, L_{e_5}^\dagger]|_{V_c}$ & $\lambda_8/2$ \\
  \bottomrule
\end{tabular}
\end{center}
where $L_x$ (resp.\ $R_x$) denotes left (resp.\ right) multiplication by $x$
in $\mathbb{O}$.

One verifies $[T_a, T_b] = if_{abc}T_c$ with the standard $SU(3)$
structure constants $f_{abc}$.

%──────────────────────────────────────────────────────────────────────────────
\section{Assessment for UBT}
\label{sec:assessment}
%──────────────────────────────────────────────────────────────────────────────

\begin{gapbox}
\textbf{Assessment: SKETCH with important caveat.}

\textbf{What works:}
\begin{itemize}
  \item The mathematical identity $SU(3) = \mathrm{Stab}_{G_2}(e_7)$ is a
        theorem in the theory of exceptional Lie groups (Theorem~\ref{thm:su3stab}).
  \item The 8 generators are explicitly identified with Gell-Mann matrices
        (Sec.~\ref{sec:gellmann}).
  \item The action on $\mathbb{C}^3$ reproduces the fundamental and adjoint
        representations.
\end{itemize}

\textbf{What does not work for UBT:}
\begin{enumerate}
  \item \textbf{UBT foundation is $\mathbb{C}\otimes\mathbb{H}$, not $\mathbb{C}\otimes\mathbb{O}$.}
    The octonions are \emph{non-associative}.  UBT requires associativity
    (the FPE, gauge covariant derivative, and action $S[\Theta]$ all use
    associative multiplication).  Octonions are used here as an auxiliary
    tool only.

  \item \textbf{Dynamical content requires additional argument.}
    The octonion approach identifies the symmetry group but does not explain
    \emph{why} the $SU(3)$ gauge symmetry is a dynamical symmetry of the UBT
    action $S[\Theta]$.  Gauge invariance of $S[\Theta]$ under $SU(3)$
    requires proof that the UBT Lagrangian is invariant under the
    $V_c$-automorphisms --- this is proved in
    \texttt{appendix\_G\_internal\_color\_symmetry.tex} via involutions,
    not octonion automorphisms.

  \item \textbf{Choice of $e_7$.}
    The identification $SU(3) = \mathrm{Stab}_{G_2}(e_7)$ breaks the
    octonionic symmetry by fixing one imaginary unit.  In the involution
    approach, the analogous symmetry breaking is performed by the $P_2$
    involution (quaternionic conjugation), which is an intrinsic operation
    within $\mathbb{C}\otimes\mathbb{H}$ --- no external element $e_7$
    is needed.
\end{enumerate}

\textbf{Conclusion:} The octonion approach establishes $SU(3)$ abstractly
but is not the proof used in UBT.  The involution approach
(see \texttt{step1\_involution\_summary.tex} and
\texttt{appendix\_G\_internal\_color\_symmetry.tex}) gives the
\textbf{proved [L0]} result within the $\mathbb{C}\otimes\mathbb{H}$
framework.
\end{gapbox}

\ifdefined\INCLUDEMODE
\else
  \end{document}
\fi
